\begin{table}[htbp]
\centering
\small
\caption{Comparação pareada, janela a janela, do DBSCAN com cada método de referência (teste de Wilcoxon).}
\label{tab:pareada}
\begin{adjustbox}{max width=\linewidth}
\begin{tabular}{lllrrrr}
\toprule
Métrica & Período & Referência & Janelas & DBSCAN melhor (\%) & $\Delta$ mediano & $p$ \\
\midrule
PAI & Treino & KDE (mesma área) & 262 & 22,5 & $-$0,11 & $<$0,001 \\
 &  & Ranking de bairros (incidência) & 262 & 83,2 & 1,32 & $<$0,001 \\
 &  & População (mesma área) & 262 & 67,6 & 1,24 & $<$0,001 \\
 &  & Todos os casos (sem clusterizar) & 262 & 88,2 & 1,65 & $<$0,001 \\
\addlinespace
 & Teste & KDE (mesma área) & 277 & 18,8 & $-$0,36 & $<$0,001 \\
 &  & Ranking de bairros (incidência) & 277 & 91,3 & 2,40 & $<$0,001 \\
 &  & População (mesma área) & 277 & 72,6 & 1,34 & $<$0,001 \\
 &  & Todos os casos (sem clusterizar) & 277 & 88,8 & 1,49 & $<$0,001 \\
\addlinespace
\midrule
PAI$_\text{pop}$ & Treino & KDE (mesma área) & 262 & 26,7 & $-$0,07 & $<$0,001 \\
 &  & Ranking de bairros (incidência) & 262 & 44,3 & 0,00 & 0,176 \\
 &  & População (mesma área) & 262 & 90,5 & 1,31 & $<$0,001 \\
 &  & Todos os casos (sem clusterizar) & 262 & 84,7 & 0,58 & $<$0,001 \\
\addlinespace
 & Teste & KDE (mesma área) & 277 & 28,2 & $-$0,10 & $<$0,001 \\
 &  & Ranking de bairros (incidência) & 277 & 45,1 & $-$0,07 & 0,286 \\
 &  & População (mesma área) & 277 & 96,8 & 1,41 & $<$0,001 \\
 &  & Todos os casos (sem clusterizar) & 277 & 83,8 & 0,55 & $<$0,001 \\
\bottomrule
\end{tabular}
\end{adjustbox}
\par\smallskip\begin{minipage}{0.95\linewidth}\footnotesize $\Delta$ = métrica do DBSCAN $-$ métrica da referência, por janela; valores positivos favorecem o DBSCAN. As janelas semanais se sobrepõem (3 das 4 semanas), de modo que os valores de $p$ são aproximados.\end{minipage}
\end{table}
